% ============================================================
\chapter{Dualit\'e en Programmation Lin\'eaire}
\label{ch:dualite}
% ============================================================

Tout probl\`eme d'optimisation cache un miroir~: son \emph{dual}. Cette id\'ee, qui remonte \`a John von Neumann et \`a ses travaux sur la th\'eorie des jeux dans les ann\'ees 1940, est l'une des plus profondes de toute l'optimisation. En programmation lin\'eaire, la dualit\'e prend une forme particuli\`erement \'el\'egante~: \`a tout probl\`eme de minimisation correspond un probl\`eme de maximisation dont la valeur optimale co\"incide avec celle du primal (th\'eor\`eme de dualit\'e forte). Le dual fournit non seulement des bornes certifi\'ees, mais aussi une interpr\'etation \'economique~: les variables duales sont des \og prix fictifs \fg qui mesurent la valeur marginale de chaque contrainte.

\section{Motivation}

Consid\'erons un probl\`eme de minimisation.  On cherche \`a encadrer la valeur
optimale~:
\begin{itemize}
  \item Toute solution r\'ealisable fournit une \textbf{borne sup\'erieure}.
  \item Peut-on obtenir syst\'ematiquement des \textbf{bornes inf\'erieures}~?
\end{itemize}
La th\'eorie de la dualit\'e r\'epond positivement \`a cette question en
associant \`a chaque PL (le \textbf{primal}) un autre PL (le \textbf{dual}).

\section{Construction du dual}

\begin{definition}[Primal et dual]
\label{def:primal_dual}
Soit le primal (P) sous forme canonique~:
\[
\text{(P)} \quad
\begin{aligned}
  \max \quad & \mathbf{c}^\top \mathbf{x} \\
  \text{s.c.} \quad & A\mathbf{x} \le \mathbf{b} \\
  & \mathbf{x} \ge \mathbf{0}
\end{aligned}
\]
Le \textbf{dual} (D) est~:
\[
\text{(D)} \quad
\begin{aligned}
  \min \quad & \mathbf{b}^\top \mathbf{y} \\
  \text{s.c.} \quad & A^\top \mathbf{y} \ge \mathbf{c} \\
  & \mathbf{y} \ge \mathbf{0}
\end{aligned}
\]
o\`u $\mathbf{y} \in \R^m$ est le vecteur des \textbf{variables duales}.
\end{definition}

\begin{formules}[title=\textbf{R\`egles de construction du dual}]
\begin{center}
\renewcommand{\arraystretch}{1.3}
\begin{tabular}{lll}
  \toprule
  \textbf{Primal (max)} & & \textbf{Dual (min)} \\
  \midrule
  Contrainte $\le$ & $\leftrightarrow$ & Variable $y_i \ge 0$ \\
  Contrainte $\ge$ & $\leftrightarrow$ & Variable $y_i \le 0$ \\
  Contrainte $=$ & $\leftrightarrow$ & Variable $y_i$ libre \\
  Variable $x_j \ge 0$ & $\leftrightarrow$ & Contrainte $\ge$ \\
  Variable $x_j \le 0$ & $\leftrightarrow$ & Contrainte $\le$ \\
  Variable $x_j$ libre & $\leftrightarrow$ & Contrainte $=$ \\
  \bottomrule
\end{tabular}
\end{center}
\end{formules}

\begin{remarque}
Le dual du dual est le primal~: $(D(D)) = (P)$.
\end{remarque}

\section{Th\'eor\`emes de dualit\'e}

\begin{theoreme}[Dualit\'e faible]
\label{thm:dualite_faible}
Si $\mathbf{x}$ est r\'ealisable pour (P) et $\mathbf{y}$ est r\'ealisable
pour (D), alors~:
\[
  \mathbf{c}^\top \mathbf{x} \le \mathbf{b}^\top \mathbf{y}
\]
\end{theoreme}

\begin{proof}
On a $A\mathbf{x} \le \mathbf{b}$ et $\mathbf{y} \ge \mathbf{0}$, donc
$\mathbf{y}^\top A\mathbf{x} \le \mathbf{y}^\top \mathbf{b}$.  De m\^eme,
$A^\top\mathbf{y} \ge \mathbf{c}$ et $\mathbf{x} \ge \mathbf{0}$ donnent
$\mathbf{x}^\top A^\top\mathbf{y} \ge \mathbf{c}^\top\mathbf{x}$.  Donc
$\mathbf{c}^\top\mathbf{x} \le \mathbf{y}^\top A\mathbf{x} \le
\mathbf{b}^\top\mathbf{y}$.
\end{proof}

\begin{corollaire}[Bornes]
\begin{enumerate}
  \item Si (P) est non born\'e ($\max = +\infty$), alors (D) est non
        r\'ealisable.
  \item Si (D) est non born\'e ($\min = -\infty$), alors (P) est non
        r\'ealisable.
\end{enumerate}
\end{corollaire}

\begin{theoreme}[Dualit\'e forte]
\label{thm:dualite_forte}
Si (P) admet une solution optimale finie $\mathbf{x}^*$, alors (D) admet
\'egalement une solution optimale finie $\mathbf{y}^*$ et~:
\[
  \mathbf{c}^\top \mathbf{x}^* = \mathbf{b}^\top \mathbf{y}^*
\]
\end{theoreme}

\begin{intuition}[title=\textbf{Interpr\'etation}]
Le th\'eor\`eme de dualit\'e forte affirme que la meilleure borne inf\'erieure
(venant du dual) co\"incide avec la meilleure borne sup\'erieure (venant du
primal) \`a l'optimum.  Il n'y a \og pas d'\'ecart \fg entre les deux.
\end{intuition}

\section{\'Ecarts compl\'ementaires}

\begin{theoreme}[\'Ecarts compl\'ementaires]
\label{thm:ecarts_complementaires}
Soient $\mathbf{x}^*$ et $\mathbf{y}^*$ des solutions r\'ealisables du primal
et du dual respectivement.  Elles sont toutes deux optimales si et seulement
si~:
\[
  y_i^* \left(b_i - \sum_{j} a_{ij} x_j^*\right) = 0, \quad i = 1, \ldots, m
\]
\[
  x_j^* \left(\sum_{i} a_{ij} y_i^* - c_j\right) = 0, \quad j = 1, \ldots, n
\]
\end{theoreme}

En d'autres termes~:
\begin{itemize}
  \item Si une contrainte primale n'est \textbf{pas satur\'ee}, la variable
        duale correspondante est nulle.
  \item Si une variable primale est \textbf{strictement positive}, la
        contrainte duale correspondante est satur\'ee.
\end{itemize}

\section{Interpr\'etation \'economique}

\begin{definition}[Prix fictif (shadow price)]
La variable duale $y_i^*$ associ\'ee \`a la $i$-\`eme contrainte du primal
s'appelle le \textbf{prix fictif} de la ressource $i$.  Elle repr\'esente la
variation marginale de la valeur optimale lorsque le membre droit $b_i$
augmente d'une unit\'e~:
\[
  y_i^* = \frac{\partial z^*}{\partial b_i}
\]
\end{definition}

\begin{exemple}[Interpr\'etation des prix fictifs]
Reprenons le probl\`eme de production du chapitre pr\'ec\'edent~:
\[
\begin{aligned}
  \max \quad & z = 5x_1 + 4x_2 \\
  \text{s.c.} \quad
    & 6x_1 + 4x_2 \le 24 \quad (y_1) \\
    & x_1 + 2x_2 \le 6 \quad (y_2) \\
    & x_1, x_2 \ge 0
\end{aligned}
\]

Le dual est~:
\[
\begin{aligned}
  \min \quad & w = 24y_1 + 6y_2 \\
  \text{s.c.} \quad
    & 6y_1 + y_2 \ge 5 \\
    & 4y_1 + 2y_2 \ge 4 \\
    & y_1, y_2 \ge 0
\end{aligned}
\]

Du tableau final du simplexe (chapitre~3), on lit~:
$y_1^* = 3/4$, $y_2^* = 1/2$.  On v\'erifie~: $w^* = 24(3/4) + 6(1/2) = 21
= z^*$.

Interpr\'etation~: augmenter la capacit\'e de la contrainte~1 de 24 \`a 25
augmente $z^*$ d'environ $3/4 = 0.75$.
\end{exemple}

\section{Lecture du dual dans le tableau du simplexe}

\begin{proposition}[Variables duales et tableau final]
Dans le tableau final du simplexe (primal en forme standard avec variables
d'\'ecart), les valeurs des variables duales optimales se lisent dans la
\textbf{ligne objectif} aux colonnes des variables d'\'ecart~:
\[
  y_i^* = \bar{c}_{s_i} \quad \text{(co\^ut r\'eduit de la variable d'\'ecart $s_i$)}
\]
\end{proposition}

\section{Le dual dans le simplexe}

La m\'ethode du \textbf{simplexe dual} r\'esout le dual directement sur le
tableau du primal.  Elle est utile quand la solution de base est
\textbf{dual-r\'ealisable} (tous les co\^uts r\'eduits $\ge 0$) mais pas
primal-r\'ealisable ($\bar{b}_i < 0$ pour certains $i$).

\begin{algorithme}[title=\textbf{Simplexe dual -- une it\'eration}]
\begin{enumerate}
  \item \textbf{Test d'optimalit\'e primal}~: si $\bar{b}_i \ge 0$ pour
        tout $i$, \textsc{stop} (solution optimale).
  \item \textbf{Variable sortante}~: choisir $r = \argmin_i \bar{b}_i$
        (la ligne la plus n\'egative).
  \item \textbf{Test de non-r\'ealisabilit\'e}~: si $\bar{a}_{rj} \ge 0$
        pour tout $j$, \textsc{stop} (primal non r\'ealisable).
  \item \textbf{Variable entrante}~: appliquer le ratio dual~:
        \[
          k = \argmin_{j : \bar{a}_{rj} < 0}
            \frac{\bar{c}_j}{\abs{\bar{a}_{rj}}}
        \]
  \item \textbf{Pivot}~: pivoter sur $\bar{a}_{rk}$.
\end{enumerate}
\end{algorithme}

\section{Exemple num\'erique du dual}

\begin{exemple}[Construction et r\'esolution du dual]
\[
\text{(P)} \quad
\begin{aligned}
  \max \quad & z = 4x_1 + 3x_2 \\
  \text{s.c.} \quad
    & 2x_1 + x_2 \le 10 \\
    & x_1 + 3x_2 \le 12 \\
    & x_1, x_2 \ge 0
\end{aligned}
\]

Le dual~:
\[
\text{(D)} \quad
\begin{aligned}
  \min \quad & w = 10y_1 + 12y_2 \\
  \text{s.c.} \quad
    & 2y_1 + y_2 \ge 4 \\
    & y_1 + 3y_2 \ge 3 \\
    & y_1, y_2 \ge 0
\end{aligned}
\]
\end{exemple}

\begin{codeblock}[title=\textbf{V\'erification avec PuLP}]
\begin{minted}{python}
import pulp

# Primal
P = pulp.LpProblem("Primal", pulp.LpMaximize)
x1 = pulp.LpVariable("x1", lowBound=0)
x2 = pulp.LpVariable("x2", lowBound=0)
P += 4*x1 + 3*x2
c1 = P.addConstraint(2*x1 + x2 <= 10, "c1")
c2 = P.addConstraint(x1 + 3*x2 <= 12, "c2")
P.solve(pulp.PULP_CBC_CMD(msg=0))
print(f"Primal : z* = {pulp.value(P.objective):.2f}")
print(f"  x* = ({x1.varValue}, {x2.varValue})")

# Dual
D = pulp.LpProblem("Dual", pulp.LpMinimize)
y1 = pulp.LpVariable("y1", lowBound=0)
y2 = pulp.LpVariable("y2", lowBound=0)
D += 10*y1 + 12*y2
D += 2*y1 + y2 >= 4
D += y1 + 3*y2 >= 3
D.solve(pulp.PULP_CBC_CMD(msg=0))
print(f"Dual : w* = {pulp.value(D.objective):.2f}")
print(f"  y* = ({y1.varValue}, {y2.varValue})")
\end{minted}
\end{codeblock}

\begin{codeoutput}
Primal : z* = 22.80
  x* = (5.4, 2.2)
Dual : w* = 22.80
  y* = (1.8, 0.4)
\end{codeoutput}

\section{R\'esum\'e des relations primal-dual}

\begin{center}
\renewcommand{\arraystretch}{1.3}
\begin{tabular}{lccc}
  \toprule
  & \textbf{(P) r\'ealisable} & \textbf{(P) non r\'ealisable} \\
  \midrule
  \textbf{(D) r\'ealisable} & Optima finis \'egaux & (D) non born\'e \\
  \textbf{(D) non r\'ealisable} & (P) non born\'e & Les deux non r\'ealisables \\
  \bottomrule
\end{tabular}
\end{center}

\begin{remarque}
Il est possible que le primal et le dual soient \emph{tous deux} non
r\'ealisables.  Exemple~: $\max\ x_1 - x_2$ s.c.\ $-x_1 + x_2 \le -1$,
$x_1 - x_2 \le -1$, $x_1, x_2 \ge 0$.
\end{remarque}

\section{Exercices}

\begin{exercice}[$\star$ -- \'Ecrire le dual]
\'Ecrire le dual des PL suivants~:
\begin{enumerate}
  \item $\max\ 3x_1 + 2x_2$ s.c.\ $x_1 + x_2 \le 5$,
        $2x_1 + x_2 \le 8$, $x_1, x_2 \ge 0$.
  \item $\min\ 5x_1 + 3x_2$ s.c.\ $x_1 + x_2 = 4$,
        $2x_1 + x_2 \ge 6$, $x_1, x_2 \ge 0$.
\end{enumerate}
\end{exercice}

\begin{exercice}[$\star\star$ -- V\'erification de la dualit\'e forte]
R\'esoudre le primal et le dual de l'exercice~1(a) et v\'erifier que les
valeurs optimales sont \'egales.
\end{exercice}

\begin{exercice}[$\star\star$ -- \'Ecarts compl\'ementaires]
Utiliser les \'ecarts compl\'ementaires pour trouver la solution duale optimale
\`a partir de la solution primale $x_1^* = 3$, $x_2^* = 2$ du PL~:
$\max\ 4x_1 + 3x_2$ s.c.\ $2x_1 + x_2 \le 8$, $x_1 + 2x_2 \le 7$,
$x_1, x_2 \ge 0$.
\end{exercice}

\begin{exercice}[$\star\star$ -- Interpr\'etation \'economique]
Une usine a trois ressources $R_1, R_2, R_3$ avec les capacit\'es 100, 80 et
120.  Les prix fictifs optimaux sont $y_1^* = 2$, $y_2^* = 0$, $y_3^* = 3$.
Interpr\'eter ces valeurs.  Un fournisseur propose 10 unit\'es
suppl\'ementaires de $R_2$ pour 5\,EUR.  Faut-il accepter~?
\end{exercice}

\begin{exercice}[$\star\star\star$ -- D\'emonstration]
D\'emontrer le th\'eor\`eme de dualit\'e faible pour le cas g\'en\'eral
(forme standard du primal).
\end{exercice}

\begin{exercice}[$\star\star\star$ -- Simplexe dual]
Appliquer le simplexe dual au probl\`eme~:
\[
\begin{aligned}
  \min \quad & z = 2x_1 + 3x_2 + x_3 \\
  \text{s.c.} \quad
    & x_1 + x_2 + x_3 \ge 6 \\
    & 2x_1 + x_3 \ge 4 \\
    & x_1, x_2, x_3 \ge 0
\end{aligned}
\]
\end{exercice}
